% Einheit 2 von 3 – Kegel (bank/pyramide-kegel-kugel/e2.jsonl, zusammenbau v0.1)
%% TODO Zweigzeile Teil 1: was hier gelernt wird (Satz in Schülersprache aus dem Einheitstitel) – Einheit 2
\einheitenkopf[e2]{Einheit 2 von 3 · Kegel}
\zweigzeile{ab Kl. 8, je nach Buch bis Kl. 10 · P10}
\setcounter{aufgabe}{17}

%% TODO Titel: Ich-kann-Satz fehlt in der Bank (Platzhalter: Kettenname) – Nr. 18
%% TODO Anweisung: ein Satz über der ersten Teilaufgabe fehlt in der Bank – Nr. 18
\begin{aufgabe}{Radius oder Durchmesser?}
\begin{teile}
\teil Eine Eistüte ist oben $6$ cm breit. Radius oder Durchmesser? Gib die andere Größe an. \leerfeld
\end{teile}
\end{aufgabe}

%% TODO Titel: Ich-kann-Satz fehlt in der Bank (Platzhalter: Kettenname) – Nr. 19
%% TODO Anweisung: ein Satz über der ersten Teilaufgabe fehlt in der Bank – Nr. 19
\begin{aufgabe}{Kegel}
%% TODO \rechenplatz{n}: Schrittzahl des Grundfalls fehlt in der Bank (nur bei fester Schreibform, 2.3 b)
\begin{teile}
\teil Die Strecke mit ? am Kegel – Höhe, Mantellinie oder Radius? Kreuze an.\\ \kreuz{Höhe}\\ \kreuz{Mantellinie}\\ \kreuz{Radius}

\kegel{1.5}{3}{}{?}{}
\teil Kegel: Radius $3$ cm, Höhe $7$ cm – Volumen? Runde auf eine Stelle. $V$ ≈ \leerfeld[cm³]
\teil Kegel: Durchmesser $8$ cm, Höhe $9$ cm – Volumen? Runde auf eine Stelle. $V$ ≈ \leerfeld[cm³]
\teil Kegel: Radius $2{,}5$ cm, Höhe $6{,}4$ cm – Volumen? Runde auf eine Stelle. $V$ ≈ \leerfeld[cm³]
\teil Ein Trichter ist innen ein Kegel mit dem Radius $6$ cm und der Höhe $11$ cm. Wie viel Milliliter fasst er? \leerfeld[ml]
\teil Kegel: Radius $16$ cm, Höhe $30$ cm – Mantellinie? $s$ = \leerfeld[cm]
\teil Kegel: Mantellinie $29$ m, Radius $20$ m – Höhe? $h$ = \leerfeld[m]
\teil Kegel: Radius $4$ cm, Mantellinie $11$ cm – Mantel? Runde auf eine Stelle. $M$ ≈ \leerfeld[cm²]
\teil Kegel: Radius $3$ cm, Mantellinie $7$ cm – Oberfläche? Runde auf eine Stelle. $O$ ≈ \leerfeld[cm²]
\teil Skizziere das Schrägbild eines Kegels mit dem Radius $2$ cm und der Höhe $5$ cm. Zeichne die Höhe gestrichelt ein und beschrifte Radius und Höhe.

\rechenplatz{5}
\teil Ein Kegel hat den Radius $3$ cm und die Mantellinie $8$ cm. Sein Netz besteht aus einem Kreis und einem Kreisausschnitt. Welchen Radius hat der Kreis, welchen der Kreisausschnitt? \leerfeld[cm] ; \leerfeld[cm]
\teil Kegel: Volumen $150$ cm³, Radius $4$ cm – Höhe? Runde auf eine Stelle. $h$ ≈ \leerfeld[cm]
\teil Ein Turm besteht aus einem $18$ m hohen Zylinder mit einem Kegeldach: Radius $3{,}5$ m, Mantellinie $6{,}2$ m. Das Dach bekommt Schiefer für $45$ € je m². Was kostet das, und wie hoch ist das Bauwerk insgesamt \hfill (P10 2026 FOR)? \leerfeld[€] ; \leerfeld[m]
\teil Ein Kegel und ein Zylinder haben beide den Radius $3$ cm und die Höhe $5$ cm. Berechne beide Volumen. Wie oft passt der Kegel in den Zylinder?
\teil Ein Kegel und ein Zylinder sind gleich hoch. Mia behauptet: „Ist der Radius des Kegels sechsmal so groß wie der Radius des Zylinders, hat der Kegel das doppelte Volumen.“ Prüfe die Aussage mit einer Rechnung \hfill (P10 2026 FOR).
\end{teile}
\end{aufgabe}

%% TODO Titel: Ich-kann-Satz fehlt in der Bank (Platzhalter: Kettenname) – Nr. 20
%% TODO Anweisung: ein Satz über der ersten Teilaufgabe fehlt in der Bank – Nr. 20
\begin{aufgabe}{Kegel – Fehler finden, Begründen, Darstellungswechsel, Anwendung}
\begin{teile}
\teil Kegel: Radius $4$ cm, Höhe $15$ cm. Paul rechnet das Volumen: \rechnung{V &= \pi \cdot 4^2 \cdot 15 \\ V &\approx 754{,}0 \text{ cm}^3} Wo ist der Fehler?
\teil Ein Kegel und ein Zylinder haben gleichen Radius und gleiche Höhe. Begründe, warum man den Kegel dreimal füllen muss, bis der Zylinder voll ist.
\teil Eine Waffeltüte ist oben $5$ cm breit und $11$ cm tief. Skizziere sie als Kegel mit der Spitze unten und beschrifte Radius und Höhe.

\rechenplatz[halb]{4}
\teil Auf dem Bauhof liegt ein kegelförmiger Haufen Streusplitt: Durchmesser $6$ m, Höhe $2$ m. Ein Lkw kann $15$ m³ laden. Reicht eine Fuhre, um den Haufen abzufahren?
\end{teile}
\end{aufgabe}
